% No preamble. Handoff catalogue, unique prefix hnc-. 17 October 2026.
% Requires hyperref (or url) for \path. No external cross-references.
\section{Annotated catalogue of the authoritative notebooks}
\label{hnc-catalogue}
This catalogue covers all thirty-nine current notebooks, including
the normalized replacement theorem completed as the final research step.
Every path is relative to the original workspace root. The first
thirty-six sources are flat files in \path{outputs/}; Notebooks 37--39
are dedicated projects with a \path{main.tex}. Those project directories
must be included even when a default search tool ignores them.

The entries summarize the complete abstracts where present, section
and theorem structure, relevant checkpoints, and saved proof/audit
history. They do not re-certify every proof. Authoritative means the
current working statement to read instead of an older draft; internal
audits and successful typesetting are not external mathematical
certification. The original source repository is separately pinned at
\texttt{e392ef45013a581664bab76b90867170b30fe10e}.

The local homogeneous leading drift and tuning chain reaches its
stated conclusion in Notebook 37. Notebooks 38--39 add actual positive
bare-measure results. The complete mixed-history, observable-source,
continuum, reconstruction and Yang--Mills identification interfaces
remain. The global status is recorded in
\path{work/final_push_existence_dependency_audit.txt}.
The companion machine-readable inventory is
\path{work/handoff_notebook_catalogue.json}; substantive corrections
and rejected shortcuts are indexed in
\path{work/handoff_corrections_register.txt}.

Dependencies below list the main upstream notebooks, rather than
claiming a complete theorem-level dependency graph. A chronological
remaining problem in an early abstract may have been settled by a
later notebook; each entry points out the relevant later development.
All detailed numerical hypotheses remain in the source theorem.


\subsection{Notebook 01: proof obligations and verified reductions}
\label{hnc-01}
\noindent\path{outputs/ym4-construction-notebook.tex}\par

\paragraph{Main results.}
\begin{itemize}
\item A master proof-obligation map for constructing a nonzero Euclidean observable sector, followed by the stronger Yang--Mills identification problem.
\item Conditional all-order factorial estimates on joint flat tests from translated fixed-side grids; a rotation criterion with angular exponential type strictly below four; and quantitative two-scale/free-kernel reductions for positive-energy nontriviality.
\item Records source-domain and running-coupling extraction interfaces and explicit counterexamples to shortcut arguments. The opening map has been updated through Notebook 39.
\end{itemize}

\paragraph{Assumptions and dependencies.}
\begin{itemize}
\item The principal original target is the homogeneous SU(2) Wilson family and separated gauge-invariant observables.
\item The all-order estimate still requires uniform translated-grid source bounds. Rotation requires the stated entire angular growth, positive imaginary-angle kernels and detection hypotheses; power counting alone does not establish them.
\item The nontriviality reduction requires actual interacting kernel accuracy and the common continuum hierarchy.
\end{itemize}

This is a foundational reduction or an independent fixed-coupling module; no earlier notebook is listed as a required upstream input.\par

\paragraph{Limits and later developments.}
\begin{itemize}
\item This is a map and collection of scoped reductions, not a completed construction or a certification of every companion proof.
\item A scalar observable sector alone does not identify the full Yang--Mills field theory or establish a mass gap.
\end{itemize}

The chronological open-problem statements in early sections must be read with the updated opening map and the latest checkpoint. Notebooks 37--38 close local drift and add genuine bare-measure results, while the full mixed-history and continuum interfaces remain.

\paragraph{Proof modules and source maps.}
\begingroup\small
\begin{itemize}
\item \path{work/rotation_research.md}
\item \path{work/compactness_research.md}
\item \path{work/all_order_research.md}
\item \path{work/source_stability_research.md}
\item \path{work/running_coupling_research.md}
\item \path{work/nontriviality_research.md}
\end{itemize}
\endgroup

\paragraph{Selected audit records.}
\begingroup\small
\begin{itemize}
\item \path{work/final_push_existence_dependency_audit.txt}
\end{itemize}
\endgroup
Corrections/scope register: C01, C02, C03, C04.\par

\subsection{Notebook 02: preserving diagonal reflection positivity}
\label{hnc-02}
\noindent\path{outputs/ym4-reflection-notebook.tex}\par

\paragraph{Main results.}
\begin{itemize}
\item A fixed-coupling construction preserving both coordinate and diagonal reflection positivity in the same infinite-volume state from the original coordinate tori.
\item Exact diagonal feature/Haar gluing, mixed site/link boundary spaces, even Schatten inequalities, local gauge-layer removal and a full reflected-source-array comparison.
\end{itemize}

\paragraph{Assumptions and dependencies.}
\begin{itemize}
\item Fixed positive coupling and compatible coordinate tori of side 2 L\textasciicircum{}M, with L odd; the finite graph, cuts and source supports must satisfy the proved placement hypotheses.
\item The transfer comparison uses full-array Jensen normalization and a volume limit, not uniqueness of Gibbs states.
\end{itemize}

This is a foundational reduction or an independent fixed-coupling module; no earlier notebook is listed as a required upstream input.\par

\paragraph{Limits and later developments.}
\begin{itemize}
\item No uniform ultraviolet estimate, rotational invariance, continuum limit or uniqueness follows from this fixed-coupling result.
\item An admissible geometric placement is not by itself a fold theorem; the feature-space gluing supplies the additional identity.
\end{itemize}

The odd-width placement and exact diagonal slab normalization supersede tentative two-frame divisibility objections in the early scratch work.

\paragraph{Proof modules and source maps.}
\begingroup\small
\begin{itemize}
\item \path{work/two_frame_geometry_research.md}
\item \path{work/diagonal_fold_research.md}
\item \path{work/gauge_layer_research.md}
\end{itemize}
\endgroup

\paragraph{Selected audit records.}
\begingroup\small
\begin{itemize}
\item \path{work/reflection_assembly_audit.md}
\item \path{work/folding_audit.md}
\end{itemize}
\endgroup
Corrections/scope register: C05.\par

\subsection{Notebook 03: current commutators and Gaussian cutoff clusters}
\label{hnc-03}
\noindent\path{outputs/ym4-current-cutoff-notebook.tex}\par

\paragraph{Main results.}
\begin{itemize}
\item An explicit covariant cutoff-current commutator estimate using first differences.
\item A convergent sharp-cutoff Gaussian cluster expansion with a zero-free logarithm and anchored activity bounds for complex covariances.
\item A direct two-source Hessian bound and thermodynamic-limit criterion for the cutoff logarithm, potentially avoiding a full strict-support expansion for a quadratic coefficient.
\end{itemize}

\paragraph{Assumptions and dependencies.}
\begin{itemize}
\item Uniform complex covariance accretivity, operator bounds, exponential decay and a fixed finite-dimensional cell description.
\item The Hessian theorem requires the actual covariance responses in fixed cutoff coordinates, including the stated localized first and second response bounds.
\end{itemize}

Selected upstream notebooks: 01.\par

\paragraph{Limits and later developments.}
\begin{itemize}
\item The abstract Gaussian theorem does not identify the actual Yang--Mills source coefficient without coordinate/domain matching.
\item Local precision in a changing gauge is not automatically the required covariance-response hypothesis.
\end{itemize}

Notebook 04 supplies strict background support under stronger hypotheses; Notebooks 18--21 provide interacting and physical-coordinate matches in the homogeneous small-field class.

\paragraph{Proof modules and source maps.}
\begingroup\small
\begin{itemize}
\item \path{work/gaussian_cutoff_cluster_research.md}
\item \path{work/gaussian_cutoff_hessian_research.md}
\item \path{work/cutoff_source_match_research.md}
\end{itemize}
\endgroup
Corrections/scope register: C04, C06.\par

\subsection{Notebook 04: strictly local Gaussian cutoff activities}
\label{hnc-04}
\noindent\path{outputs/ym4-gaussian-locality-notebook.tex}\par

\paragraph{Main results.}
\begin{itemize}
\item A support-valued density expansion yields strictly local holomorphic Gaussian cutoff activities.
\item The proof preserves marked-set compatibility while separately tracking background supports, inverse marginal matrices and re-rooting costs.
\end{itemize}

\paragraph{Assumptions and dependencies.}
\begin{itemize}
\item A background-independent uniformly accretive reference covariance with weighted Schur bounds.
\item A small symmetric local precision perturbation on a product complex background domain, with the explicit weighted support norm.
\end{itemize}

Selected upstream notebooks: 03.\par

\paragraph{Limits and later developments.}
\begin{itemize}
\item The conclusion concerns locality, holomorphy and norms; it does not assert that each activity is gauge invariant or hypercubic invariant.
\item Uniform estimates about arbitrary reference backgrounds do not imply strict locality in those reference backgrounds.
\end{itemize}

The title and theorem wording retain some candidate-era language. Its internally audited abstract theorem is used later with actual homogeneous precision decompositions; Notebook 18 extends it to overlapping interactions.

\paragraph{Proof modules and source maps.}
\begingroup\small
\begin{itemize}
\item \path{work/gaussian_cutoff_locality_research.md}
\end{itemize}
\endgroup

\paragraph{Selected audit records.}
\begingroup\small
\begin{itemize}
\item \path{work/gaussian_cutoff_locality_audit.md}
\end{itemize}
\endgroup
Corrections/scope register: C06.\par

\subsection{Notebook 05: the flat gauge bridge and cutoff interactions}
\label{hnc-05}
\noindent\path{outputs/ym4-gauge-interaction-notebook.tex}\par

\paragraph{Main results.}
\begin{itemize}
\item A quantitative flat axial-to-exponential gauge bridge using the actual averaged-path gauge penalty, together with current-dependent Ward corrections.
\item A disjoint-cell interacting Gaussian cutoff theorem and a corrected physical-field radius estimate, including the linear Jacobian contribution.
\end{itemize}

\paragraph{Assumptions and dependencies.}
\begin{itemize}
\item The precise rooted/block-mean gauge transformations and averaging constraints are distinguished.
\item The initial bridge uses an axial coercivity input; cell interactions have bounded local analytic coefficients on fixed-dimensional disjoint cells.
\end{itemize}

Selected upstream notebooks: 03, 04.\par

\paragraph{Limits and later developments.}
\begin{itemize}
\item Disjoint-cell integration alone does not handle arbitrary overlapping polymer supports.
\item The physical fluctuation radius and the normalized integration-variable radius are different; current-free flat identities do not automatically extend to arbitrary effective currents.
\end{itemize}

Notebook 06 supplies direct flat coercivity, and Notebook 18 closes the overlapping-support integration problem. Subsequent physical source matching is in Notebooks 17--21.

\paragraph{Proof modules and source maps.}
\begingroup\small
\begin{itemize}
\item \path{work/exponential_gauge_bridge_research.md}
\item \path{work/interaction_cutoff_research.md}
\end{itemize}
\endgroup

\paragraph{Selected audit records.}
\begingroup\small
\begin{itemize}
\item \path{work/exponential_gauge_bridge_audit.md}
\item \path{work/interaction_cutoff_audit.md}
\end{itemize}
\endgroup
Corrections/scope register: C07, C08.\par

\subsection{Notebook 06: direct flat Maxwell bounds and a compact lift}
\label{hnc-06}
\noindent\path{outputs/ym4-flat-maxwell-notebook.tex}\par

\paragraph{Main results.}
\begin{itemize}
\item A curvature-averaging cochain identity and scaled L2 contraction give a lower Maxwell energy bound.
\item Explicit compact lifts give a uniform upper energy bound; finite-block tree/averaging Poincare inequalities give a retained axial gap.
\item A boundary-adapted lift preserves precisely the incident-bond Dirichlet convention and all zero-extended averaging rows.
\end{itemize}

\paragraph{Assumptions and dependencies.}
\begin{itemize}
\item Odd homogeneous averaging ratios and fixed elementary block geometry; periodic multiplicities are retained.
\item Matching a classical effective Hessian requires the actual C2 constrained-minimizer branch and feasible curves, not merely differentiability.
\end{itemize}

Selected upstream notebooks: 05.\par

\paragraph{Limits and later developments.}
\begin{itemize}
\item These are flat variational bounds, not nonflat covariance estimates or a quantum mass gap.
\item Freezing boundary-crossing bonds would change the boundary theorem's hypotheses.
\end{itemize}

The nonlinear minimizer and exact homogeneous Hessian match are supplied in Notebooks 15 and 17.

\paragraph{Proof modules and source maps.}
\begingroup\small
\begin{itemize}
\item \path{work/flat_maxwell_gap_research.md}
\item \path{work/flat_maxwell_lift_research.md}
\item \path{work/flat_maxwell_boundary_research.md}
\end{itemize}
\endgroup

\paragraph{Selected audit records.}
\begingroup\small
\begin{itemize}
\item \path{work/flat_maxwell_gap_audit.md}
\item \path{work/flat_maxwell_lift_audit.md}
\item \path{work/flat_maxwell_boundary_audit.md}
\end{itemize}
\endgroup
Corrections/scope register: C09, C10.\par

\subsection{Notebook 07: locality and refinement limit of the flat Gaussian}
\label{hnc-07}
\noindent\path{outputs/ym4-flat-locality-notebook.tex}\par

\paragraph{Main results.}
\begin{itemize}
\item Exact Fourier alias pole cancellation gives a common holomorphic strip and exponential locality of the flat effective precision and covariance.
\item The normalized midpoint quadratic symbol has the Maxwell coefficient one.
\item The explicit sinc-alias refinement limit has an O(M\textasciicircum{}-2) exponentially weighted kernel error, uniformly in compatible periodic volume.
\end{itemize}

\paragraph{Assumptions and dependencies.}
\begin{itemize}
\item The homogeneous flat straight-string averaging family, its precise normalization and the retained gauge coordinates.
\item Complex pole cancellation is performed before estimates; zero modes and finite-period multiplicities are treated explicitly.
\end{itemize}

Selected upstream notebooks: 05, 06.\par

\paragraph{Limits and later developments.}
\begin{itemize}
\item The limiting object is a free Gaussian reference, not an interacting continuum Yang--Mills theory.
\item Ellipticity alone would not establish the normalized Maxwell coefficient; the explicit symbol calculation does.
\end{itemize}

This is the flat reference used by the nonlinear, determinant and heat-comparison notebooks.

\paragraph{Proof modules and source maps.}
\begingroup\small
\begin{itemize}
\item \path{work/flat_covariance_limit_fragment.tex}
\item \path{work/flat_covariance_decay_research.md}
\item \path{work/flat_covariance_limit_research.md}
\item \path{work/flat_covariance_source_research.md}
\end{itemize}
\endgroup

\paragraph{Selected audit records.}
\begingroup\small
\begin{itemize}
\item \path{work/flat_covariance_decay_audit.md}
\item \path{work/flat_covariance_combined_audit.md}
\end{itemize}
\endgroup
Corrections/scope register: C10, C11.\par

\subsection{Notebook 08: a local flat minimizing extension}
\label{hnc-08}
\noindent\path{outputs/ym4-flat-minimizer-notebook.tex}\par

\paragraph{Main results.}
\begin{itemize}
\item An explicit minimizing extension in the deposited homogeneous source gauge, with exact constraint and gauge uniqueness.
\item Uniform exponential locality in physical cell norms for the potential, first derivatives, curvature and componentwise Laplacian; C1,alpha regularity for every alpha below one.
\item A fixed-constraint first curvature-metric variation is identified.
\end{itemize}

\paragraph{Assumptions and dependencies.}
\begin{itemize}
\item The source scalar projection uses the squared-Laplacian scalar correction, distinct from the least-L2 minimizing lift.
\item Flat homogeneous averaging and compatible periods, with physical versus counting normalizations explicit.
\end{itemize}

Selected upstream notebooks: 06, 07.\par

\paragraph{Limits and later developments.}
\begin{itemize}
\item No endpoint Lipschitz bound on first derivatives or general bounded individual second differences is asserted.
\item A flat local extension is not yet a nonflat/current-domain construction.
\end{itemize}

Notebooks 9--16 build the inverse, current, regularity and nonlinear consequences. Notebook 10 completes the curvature-response resolvent.

\paragraph{Proof modules and source maps.}
\begingroup\small
\begin{itemize}
\item \path{work/flat_minimizer_locality_research.md}
\item \path{work/flat_minimizer_source_gauge_research.md}
\item \path{work/flat_curvature_response_research.md}
\end{itemize}
\endgroup

\paragraph{Selected audit records.}
\begingroup\small
\begin{itemize}
\item \path{work/flat_minimizer_locality_audit.md}
\end{itemize}
\endgroup
Corrections/scope register: C10, C12.\par

\subsection{Notebook 09: the flat fine inverse and a current neighborhood}
\label{hnc-09}
\noindent\path{outputs/ym4-flat-current-notebook.tex}\par

\paragraph{Main results.}
\begin{itemize}
\item A complete Woodbury formula for the flat fine inverse with cancellation of the apparent central poles.
\item Uniform local physical-cell operator bounds with at most two total fine derivatives.
\item A small complex auxiliary-current neighborhood, localized factorial derivatives and a quadratic zero-current remainder of the coarse precision.
\end{itemize}

\paragraph{Assumptions and dependencies.}
\begin{itemize}
\item The homogeneous source operators and a positive auxiliary penalty in a fixed compact interval; the stated unit-scale convention is identified separately.
\item The current enters through the proved source-response factorization.
\end{itemize}

Selected upstream notebooks: 07, 08.\par

\paragraph{Limits and later developments.}
\begin{itemize}
\item The separate nonlinear constraint-current term is not covered by the auxiliary-current result alone.
\item An unprinted scale-dependent penalty recursion is not assumed.
\end{itemize}

Notebook 11 adds the missing constraint-current term. Notebook 12 gives pointwise bounds in the weak source norms.

\paragraph{Proof modules and source maps.}
\begingroup\small
\begin{itemize}
\item \path{work/flat_auxiliary_current_fragment.tex}
\item \path{work/flat_full_inverse_fragment.tex}
\item \path{work/flat_full_inverse_research.md}
\item \path{work/flat_scalar_gauge_locality_research.md}
\item \path{work/flat_auxiliary_current_neighborhood_research.md}
\end{itemize}
\endgroup

\paragraph{Selected audit records.}
\begingroup\small
\begin{itemize}
\item \path{work/flat_full_inverse_audit.md}
\item \path{work/flat_scalar_gauge_locality_audit.md}
\item \path{work/flat_auxiliary_current_neighborhood_audit.md}
\item \path{work/flat_current_integrated_audit.md}
\end{itemize}
\endgroup
Corrections/scope register: C12, C13.\par

\subsection{Notebook 10: local curvature perturbations at a fixed constraint}
\label{hnc-10}
\noindent\path{outputs/ym4-flat-curvature-notebook.tex}\par

\paragraph{Main results.}
\begin{itemize}
\item The constrained curvature covariance is the orthogonal projection onto the closure of feasible curls.
\item A small fixed-constraint curvature-metric perturbation has an exact resolvent formula and a weighted local expansion at all orders.
\end{itemize}

\paragraph{Assumptions and dependencies.}
\begin{itemize}
\item The averaging constraint is fixed; the perturbing metric is a bounded local operator in physical curvature cell norms.
\item The projection is defined on the closure of feasible curls, without assuming every element has an attained potential minimizer.
\end{itemize}

Selected upstream notebooks: 08, 09.\par

\paragraph{Limits and later developments.}
\begin{itemize}
\item This does not establish the actual nonlinear background/current dependence of an RG Hessian.
\item Higher response locality uses the local projection, not only a norm bound on it.
\end{itemize}

The projection supplies an independent fixed-constraint perturbative module; later nonflat results use a separate actual Wilson form calculation.

\paragraph{Proof modules and source maps.}
\begingroup\small
\begin{itemize}
\item \path{work/flat_curvature_projection_fragment.tex}
\item \path{work/flat_curvature_response_research.md}
\end{itemize}
\endgroup

\paragraph{Selected audit records.}
\begingroup\small
\begin{itemize}
\item \path{work/flat_curvature_projection_audit.md}
\end{itemize}
\endgroup

\subsection{Notebook 11: the full flat constraint-current family}
\label{hnc-11}
\noindent\path{outputs/ym4-full-flat-current-notebook.tex}\par

\paragraph{Main results.}
\begin{itemize}
\item Composition of fixed-elementary-scale quadratic average jets bounds the full nonlinear constraint-current operator in Hilbert and supremum norms.
\item A common weighted complex-current disk controls the complete flat inverse, extension and coarse precision, with factorial localized responses.
\end{itemize}

\paragraph{Assumptions and dependencies.}
\begin{itemize}
\item Homogeneous depths M=L\textasciicircum{}j with fixed L; scale factors are summed rather than treating each stage independently.
\item The current is bounded; the proof does not differentiate it.
\end{itemize}

Selected upstream notebooks: 08, 09.\par

\paragraph{Limits and later developments.}
\begin{itemize}
\item Mixed histories and nonflat current domains remain separate.
\item The complete source current cannot be replaced by the auxiliary current of Notebook 09.
\end{itemize}

Notebook 12 transfers this complete family to the actual homogeneous weak regularity norms.

\paragraph{Proof modules and source maps.}
\begingroup\small
\begin{itemize}
\item \path{work/flat_full_constraint_current_fragment.tex}
\item \path{work/flat_full_constraint_current_research.md}
\end{itemize}
\endgroup

\paragraph{Selected audit records.}
\begingroup\small
\begin{itemize}
\item \path{work/flat_full_constraint_current_audit.md}
\item \path{work/flat_full_current_integrated_audit.md}
\end{itemize}
\endgroup
Corrections/scope register: C13.\par

\subsection{Notebook 12: fine regularity on the flat current domain}
\label{hnc-12}
\noindent\path{outputs/ym4-flat-regularity-notebook.tex}\par

\paragraph{Main results.}
\begin{itemize}
\item Arbitrary bounded fine forcing is mapped by the flat inverse to C1,alpha for alpha below one, with exponential cell locality.
\item Exact factorizations transfer bounded currents to N-to-N\_beta estimates without derivatives of the current.
\item Right-divergence estimates, an equal-exponent Schauder improvement, scalar inverse analogues and the separate constrained-current member are proved.
\end{itemize}

\paragraph{Assumptions and dependencies.}
\begin{itemize}
\item Fixed exponents strictly between zero and one, bounded currents in the common Hilbert/weak disk, and complete-field weighted norms.
\item The scalar inverse retains its explicit positive penalty condition; source localization uses smooth collars or global weighted estimates.
\end{itemize}

Selected upstream notebooks: 08, 09, 11.\par

\paragraph{Limits and later developments.}
\begin{itemize}
\item The gauge projection is not asserted to be bounded on arbitrary L-infinity data.
\item The zero-current Laplacian bound is not extended to nonzero currents; endpoint and individual second-derivative bounds are excluded.
\end{itemize}

The constrained equal-exponent Green estimates are used in Notebook 15's nonlinear minimizer.

\paragraph{Proof modules and source maps.}
\begingroup\small
\begin{itemize}
\item \path{work/flat_pointwise_current_fragment.tex}
\item \path{work/flat_pointwise_inverse_fragment.tex}
\item \path{work/flat_right_divergence_fragment.tex}
\item \path{work/flat_schauder_scalar_fragment.tex}
\item \path{work/source_flat_weak_norm_match.md}
\end{itemize}
\endgroup

\paragraph{Selected audit records.}
\begingroup\small
\begin{itemize}
\item \path{work/flat_pointwise_inverse_audit.md}
\item \path{work/flat_pointwise_current_audit.md}
\item \path{work/flat_right_divergence_audit.md}
\item \path{work/flat_schauder_scalar_audit.md}
\item \path{work/flat_regularity_integrated_audit.md}
\end{itemize}
\endgroup
Corrections/scope register: C12, C13.\par

\subsection{Notebook 13: a uniform nonlinear constraint chart}
\label{hnc-13}
\noindent\path{outputs/ym4-nonlinear-constraint-notebook.tex}\par

\paragraph{Main results.}
\begin{itemize}
\item A refinement-uniform analytic ball for the complete fixed-reference corrected nonlinear average.
\item An exact nonlinear constraint chart and a uniformly local right inverse of its ordinary differential.
\end{itemize}

\paragraph{Assumptions and dependencies.}
\begin{itemize}
\item Fixed primitive blocking, small global physical fields, and the specific fixed-reference corrected averaging convention.
\item The uniform estimate uses the flat contraction together with the summable geometric scale factors.
\end{itemize}

Selected upstream notebooks: 08, 12.\par

\paragraph{Limits and later developments.}
\begin{itemize}
\item The derivative in a moving reference is a different operator and requires an endpoint correction.
\item This is a constraint chart, not yet the nonlinear minimum.
\end{itemize}

Notebook 14 identifies the moving-reference relation; Notebook 15 constructs the minimum.

\paragraph{Proof modules and source maps.}
\begingroup\small
\begin{itemize}
\item \path{work/flat_nonlinear_average_fragment.tex}
\item \path{work/flat_nonlinear_average_research.md}
\end{itemize}
\endgroup

\paragraph{Selected audit records.}
\begingroup\small
\begin{itemize}
\item \path{work/flat_nonlinear_average_audit.md}
\end{itemize}
\endgroup
Corrections/scope register: C10.\par

\subsection{Notebook 14: nonflat backgrounds and energy stability}
\label{hnc-14}
\noindent\path{outputs/ym4-nonflat-background-notebook.tex}\par

\paragraph{Main results.}
\begin{itemize}
\item The exact endpoint correction relates fixed- and moving-reference average derivatives; the ordinary reference trajectory and both actual averages are controlled.
\item An explicit Wilson plaquette Hessian calculation yields a uniformly small local form perturbation on small global physical C1 backgrounds.
\item The scalar compression and moving gauge projection give a stable energy inverse; the gauge projection is independent of the auxiliary scalar penalty when defined.
\end{itemize}

\paragraph{Assumptions and dependencies.}
\begin{itemize}
\item The actual homogeneous averages and a small globally defined C1 connection, with the normalized four-dimensional bare Wilson action and left-exponential chart.
\item Only energy blocks with at most one derivative on each side are used.
\end{itemize}

Selected upstream notebooks: 09, 12, 13.\par

\paragraph{Limits and later developments.}
\begin{itemize}
\item No general nonflat pointwise/current theorem or arbitrary mixed-history domain follows.
\item The deposited normalization match is separated from broader original-paper graded-curl/effective-action conventions.
\end{itemize}

This provides actual small-background energy control for the subsequent nonlinear construction, not a substitute for that construction.

\paragraph{Proof modules and source maps.}
\begingroup\small
\begin{itemize}
\item \path{work/nonflat_averaging_coordinate_fragment.tex}
\item \path{work/nonflat_reference_averages_fragment.tex}
\item \path{work/nonflat_wilson_inverse_fragment.tex}
\item \path{work/nonflat_wilson_inverse_research.md}
\end{itemize}
\endgroup

\paragraph{Selected audit records.}
\begingroup\small
\begin{itemize}
\item \path{work/nonflat_averaging_coordinate_audit.md}
\item \path{work/nonflat_reference_averages_audit.md}
\item \path{work/nonflat_wilson_inverse_audit.md}
\item \path{work/nonflat_background_integrated_audit.md}
\end{itemize}
\endgroup
Corrections/scope register: C14.\par

\subsection{Notebook 15: a uniform nonlinear constrained Wilson minimum}
\label{hnc-15}
\noindent\path{outputs/ym4-nonlinear-minimizer-notebook.tex}\par

\paragraph{Main results.}
\begin{itemize}
\item An equal-exponent constrained Green operator, exact Wilson force and two-Hilbert constraint-curvature estimate produce a uniform analytic fixed-slice minimum.
\item A controlled constraint-preserving gauge flow covers nearby gauge orbits.
\item A separate ordinary-average transfer yields a local constrained minimum modulo gauge in the small global C1,alpha class.
\end{itemize}

\paragraph{Assumptions and dependencies.}
\begin{itemize}
\item Finite homogeneous periodic lattices, fixed primitive blocking and the uniform small global physical field domain.
\item The minimum comparison is on a common strict small ball; the gauge-transfer path remains inside the allowed domain.
\end{itemize}

Selected upstream notebooks: 12, 13, 14.\par

\paragraph{Limits and later developments.}
\begin{itemize}
\item This is a local classical result, not a global minimum among arbitrary Wilson configurations.
\item It does not supply large-field integration, arbitrary topological sectors or mixed histories.
\end{itemize}

The C2/analytic minimum closes the flat Hessian identification used earlier. Later notebooks localize its responses and scalar action.

\paragraph{Proof modules and source maps.}
\begingroup\small
\begin{itemize}
\item \path{work/flat_constrained_schauder_fragment.tex}
\item \path{work/nonlinear_fixed_gauge_minimizer_fragment.tex}
\item \path{work/ordinary_constraint_transfer_fragment.tex}
\item \path{work/wilson_nonlinear_force_fragment.tex}
\item \path{work/nonlinear_gauge_slice_research.tex}
\end{itemize}
\endgroup

\paragraph{Selected audit records.}
\begingroup\small
\begin{itemize}
\item \path{work/flat_constrained_schauder_audit.md}
\item \path{work/nonlinear_fixed_gauge_minimizer_provenance_audit.md}
\item \path{work/nonlinear_gauge_slice_second_audit.md}
\item \path{work/ordinary_constraint_transfer_second_audit.md}
\item \path{work/nonlinear_minimizer_integrated_audit.md}
\end{itemize}
\endgroup
Corrections/scope register: C09, C15.\par

\subsection{Notebook 16: local nonlinear responses and the classical coarse action}
\label{hnc-16}
\noindent\path{outputs/ym4-nonlinear-response-notebook.tex}\par

\paragraph{Main results.}
\begin{itemize}
\item One-weight analytic responses of both minimizer representatives and their multiplier are exponentially localized at every order.
\item The ordinary representative has a controlled local gauge transfer and a distinct linear lift differing from the source lift by a gradient.
\item The envelope identity identifies the classical coarse force and Hessian; a separate constrained-Gaussian calculation proves an obstruction to refinement-uniform smooth fluctuation balls.
\end{itemize}

\paragraph{Assumptions and dependencies.}
\begin{itemize}
\item The same small nonconstant homogeneous classical backgrounds and physical norms as the nonlinear minimum.
\item Weighted differentiability is at unweighted small base fields; no unsupported weighted-smallness assumption is introduced.
\end{itemize}

Selected upstream notebooks: 15.\par

\paragraph{Limits and later developments.}
\begin{itemize}
\item One weighted response direction is not a strict multiple-support expansion.
\item The smooth Gaussian ball obstruction rules out treating the actual rough fluctuations as uniformly smooth physical perturbations.
\end{itemize}

Notebook 21 supplies strict multi-support localization by a separate support-algebra fixed point; Notebook 17 treats rough fluctuation tails.

\paragraph{Proof modules and source maps.}
\begingroup\small
\begin{itemize}
\item \path{work/classical_coarse_action_fragment.tex}
\item \path{work/gaussian_smooth_ball_obstruction_fragment.tex}
\item \path{work/nonlinear_minimizer_weighted_response_fragment.tex}
\item \path{work/ordinary_representative_locality_fragment.tex}
\item \path{work/classical_coarse_action_research.md}
\item \path{work/nonlinear_minimizer_weighted_response_research.md}
\item \path{work/ordinary_representative_locality_research.md}
\end{itemize}
\endgroup

\paragraph{Selected audit records.}
\begingroup\small
\begin{itemize}
\item \path{work/nonlinear_minimizer_weighted_response_audit.md}
\item \path{work/ordinary_representative_locality_audit.md}
\item \path{work/classical_coarse_action_audit.md}
\item \path{work/gaussian_smooth_ball_obstruction_provenance_audit.md}
\end{itemize}
\endgroup
Corrections/scope register: C15, C16.\par

\subsection{Notebook 17: joint fluctuation tails and the one-step Gaussian domain}
\label{hnc-17}
\noindent\path{outputs/ym4-gaussian-step-notebook.tex}\par

\paragraph{Main results.}
\begin{itemize}
\item The actual fixed-ratio fluctuation rescaling is identified, including the constraint-curvature term in the first reduced Hessian.
\item Joint real Gaussian bad-block tails and complex local marginal bounds avoid a total-volume loss.
\item The homogeneous flat current-step covariance is matched uniformly over earlier classical depths.
\end{itemize}

\paragraph{Assumptions and dependencies.}
\begin{itemize}
\item Fixed cell size and dimension, the explicit independent retained coordinates, and uniform covariance/coercivity bounds.
\item Possible correlations and singular observable covariances are treated rather than replaced by independence.
\end{itemize}

Selected upstream notebooks: 06, 15, 16.\par

\paragraph{Limits and later developments.}
\begin{itemize}
\item These are Gaussian fluctuation estimates, not bounds on the normalized full Wilson measure.
\item The retained-window convention and its physical versus normalized thresholds must be preserved.
\end{itemize}

The actual physical complex precision and interaction matches are supplied by Notebooks 19--21.

\paragraph{Proof modules and source maps.}
\begingroup\small
\begin{itemize}
\item \path{work/current_scale_gaussian_fragment.tex}
\item \path{work/gaussian_joint_block_tail_fragment.tex}
\item \path{work/one_step_gaussian_scaling_fragment.tex}
\item \path{work/one_step_quantum_domains_source_research.md}
\end{itemize}
\endgroup

\paragraph{Selected audit records.}
\begingroup\small
\begin{itemize}
\item \path{work/one_step_gaussian_scaling_audit.md}
\item \path{work/gaussian_joint_block_tail_provenance_audit.md}
\end{itemize}
\endgroup
Corrections/scope register: C08, C17.\par

\subsection{Notebook 18: integrating overlapping cutoff interactions}
\label{hnc-18}
\noindent\path{outputs/ym4-overlapping-integration-notebook.tex}\par

\paragraph{Main results.}
\begin{itemize}
\item A convergent Gaussian cutoff expansion for overlapping hyperedge interactions, with explicit weighted analytic coefficient norms.
\item Molecule formation, cell-endpoint Cauchy estimates and planted-tree bounds control arbitrary overlaps without an uncontrolled degree-dependent support cost.
\item Strict background support, periodic geometry and a tunable real-baseline density comparison give zero-free logarithms and quantitative interaction differences.
\end{itemize}

\paragraph{Assumptions and dependencies.}
\begin{itemize}
\item Uniform Gaussian spectral/decay bounds; bounded analytic interactions on the cutoff; small weighted coefficient norms.
\item The strict-support extension requires small local symmetric precision pieces on a product background domain. The tunable refinement starts from a centered real positive Gaussian and a controlled relative complex perturbation.
\end{itemize}

Selected upstream notebooks: 03, 04, 05, 17.\par

\paragraph{Limits and later developments.}
\begin{itemize}
\item Large local dimension is paid explicitly: choosing the relative perturbation of order inverse dimension controls the density cost but worsens fixed derivative constants.
\item Abstract overlap integration alone does not verify actual RG activity norms or domains.
\end{itemize}

Notebooks 19--21 match the theorem to the homogeneous physical step; 26 and 28 resolve support-weight and common-radius transport issues.

\paragraph{Proof modules and source maps.}
\begingroup\small
\begin{itemize}
\item \path{work/overlapping_interaction_background_fragment.tex}
\item \path{work/overlapping_interaction_cutoff_fragment.tex}
\item \path{work/periodic_overlap_locality_fragment.tex}
\item \path{work/tunable_gaussian_overlap_fragment.tex}
\item \path{work/overlapping_interaction_cutoff_research.md}
\end{itemize}
\endgroup

\paragraph{Selected audit records.}
\begingroup\small
\begin{itemize}
\item \path{work/overlapping_interaction_cutoff_audit.md}
\item \path{work/overlapping_interaction_tunable_audit.md}
\end{itemize}
\endgroup
Corrections/scope register: C06, C07.\par

\subsection{Notebook 19: the physical background Gaussian in the source coordinates}
\label{hnc-19}
\noindent\path{outputs/ym4-physical-step-notebook.tex}\par

\paragraph{Main results.}
\begin{itemize}
\item The actual one-pivot-per-constraint chart is retained, together with the fixed independent retained-link cutoff.
\item The physical gauge-zero stationary background yields an exact reduced precision containing both constraint curvature and relative-chart current curvature.
\item Local pivot Jacobian, Haar and exponential-gauge remainder bounds are proved with finite supports and bounded incidence; a uniform complex covariance disk follows.
\end{itemize}

\paragraph{Assumptions and dependencies.}
\begin{itemize}
\item Fixed elementary ratio and a small product domain of current backgrounds; constants uniform in preceding homogeneous depth.
\item Equality with the physical integral is asserted on the actual stationary real family. Background-only delta determinants are kept outside the fluctuation interaction.
\end{itemize}

Selected upstream notebooks: 15, 16, 17, 18.\par

\paragraph{Limits and later developments.}
\begin{itemize}
\item An extra reconstructed-field or pivot indicator would be moving and is not silently replaced by the retained window.
\item The analytic centered extension off the physical family is not itself an equality to a stationary physical integral.
\end{itemize}

This is the coordinate and measure input used by the bare and arbitrary-depth one-step integrations.

\paragraph{Proof modules and source maps.}
\begingroup\small
\begin{itemize}
\item \path{work/physical_current_pivot_precision_fragment.tex}
\item \path{work/physical_pivot_measure_fragment.tex}
\item \path{work/physical_current_pivot_precision_research.md}
\item \path{work/physical_pivot_measure_research.md}
\item \path{work/pivot_cutoff_source_research.md}
\end{itemize}
\endgroup

\paragraph{Selected audit records.}
\begingroup\small
\begin{itemize}
\item \path{work/physical_current_pivot_precision_audit.md}
\item \path{work/physical_pivot_measure_provenance_audit.md}
\item \path{work/physical_stationary_background_audit.md}
\end{itemize}
\endgroup
Corrections/scope register: C08, C18.\par

\subsection{Notebook 20: the first Wilson small-field integration}
\label{hnc-20}
\noindent\path{outputs/ym4-bare-step-notebook.tex}\par

\paragraph{Main results.}
\begin{itemize}
\item An actual bare Wilson small-field integration with explicit local precision and cubic action remainder.
\item The fixed retained cutoff integral has a volume-uniform connected logarithm with error controlled by g times the cutoff radius cubed plus an exponential Gaussian tail.
\item The exact group gauge insertion, Haar/pivot factors and separate delta/Gaussian determinants are retained.
\end{itemize}

\paragraph{Assumptions and dependencies.}
\begin{itemize}
\item The first classical depth M=1, the physical stationary family and its independent retained coordinates.
\item A specified gauge-invariant admissible orbit domain is required for the global insertion identity; the local chart integral is stated independently.
\end{itemize}

Selected upstream notebooks: 18, 19.\par

\paragraph{Limits and later developments.}
\begin{itemize}
\item This is one admitted small-field sector, not an estimate of its complement in the full measure.
\item Strict supports are in current backgrounds before composition with the next stationary field.
\end{itemize}

Notebook 21 extends the reference action to all homogeneous classical depths. Notebook 38 supplies a positive full-measure embedding, without identifying its remainder with a controlled mixed-history norm.

\paragraph{Proof modules and source maps.}
\begingroup\small
\begin{itemize}
\item \path{work/bare_step_integration_fragment.tex}
\end{itemize}
\endgroup

\paragraph{Selected audit records.}
\begingroup\small
\begin{itemize}
\item \path{work/bare_step_integration_audit.md}
\end{itemize}
\endgroup
Corrections/scope register: C18, C19.\par

\subsection{Notebook 21: local classical actions through arbitrary refinement}
\label{hnc-21}
\noindent\path{outputs/ym4-classical-localization-notebook.tex}\par

\paragraph{Main results.}
\begin{itemize}
\item A simultaneous support-valued fixed point gives the classical minimizer, multiplier and scalar action at arbitrary homogeneous depth.
\item Arbitrary fixed cardinality reserve is available after shrinking the background radius. Differentiating scalar action pieces gives individually symmetric local precision pieces.
\item Composition with the actual fixed-ratio physical chart closes one small-field integration of the classical reference action at every old depth.
\end{itemize}

\paragraph{Assumptions and dependencies.}
\begin{itemize}
\item Concrete smooth physical-cell partitions and full mixed-insertion coefficient norms, not just separate pointwise derivatives.
\item Composition uses rooted graph Steiner-tree length; the joint field-type maximum is taken inside each coefficient norm before summing supports.
\item The ordinary point-rooted adjoint is avoided because it is not uniformly bounded in the required physical normalization.
\end{itemize}

Selected upstream notebooks: 15, 16, 18, 19, 20.\par

\paragraph{Limits and later developments.}
\begin{itemize}
\item The result does not yet transport inherited quantum activities or close repeated quantum domains.
\item Arbitrary cardinality reserve is bought by a smaller radius; it is not a free uniform-in-all-weights assertion.
\end{itemize}

Notebook 22 transports these supported activities; 26--28 close the specific weight and radius issues.

\paragraph{Proof modules and source maps.}
\begingroup\small
\begin{itemize}
\item \path{work/classical_local_vertex_bounds_fragment.tex}
\item \path{work/classical_step_integration_fragment.tex}
\item \path{work/classical_support_localization_fragment.tex}
\item \path{work/classical_local_vertex_bounds_research.md}
\item \path{work/classical_support_localization_research.md}
\end{itemize}
\endgroup

\paragraph{Selected audit records.}
\begingroup\small
\begin{itemize}
\item \path{work/classical_support_localization_audit.md}
\item \path{work/classical_support_localization_provenance_audit.md}
\end{itemize}
\endgroup
Corrections/scope register: C16, C20.\par

\subsection{Notebook 22: transport into the next coarse background}
\label{hnc-22}
\noindent\path{outputs/ym4-coarse-transport-notebook.tex}\par

\paragraph{Main results.}
\begin{itemize}
\item The actual stationary current background is constructed as a strictly supported analytic function of next-coarse data.
\item Holomorphic activity substitution preserves the original carrier and pays explicit coordinate-count, support and radius costs.
\item A connected-carrier projection bound improves the coarsening weight estimate.
\end{itemize}

\paragraph{Assumptions and dependencies.}
\begin{itemize}
\item Fixed-dimensional current-block coordinates and the actual affine pivot chart; the gradient majorant is obtained coefficient by coefficient.
\item The substitution proof uses explicit colored scalar slots where a maximum over separate input families is legitimate.
\item Input supports, current Gaussian cells and subsequent coarsening labels are not identified or projected twice.
\end{itemize}

Selected upstream notebooks: 19, 21.\par

\paragraph{Limits and later developments.}
\begin{itemize}
\item The map is bounded with explicit losses, not an invariant analytic-ball or amplitude contraction theorem.
\item The pure-gauge counterexample forbids a blanket inverse-ratio supremum contraction in the gauge-zero representative.
\end{itemize}

The smaller pre-gauge representative and auxiliary grouping in 25--28 address these particular losses.

\paragraph{Proof modules and source maps.}
\begingroup\small
\begin{itemize}
\item \path{work/coarse_stationary_support_fragment.tex}
\item \path{work/quantum_activity_transport_fragment.tex}
\item \path{work/coarse_stationary_support_research.md}
\item \path{work/quantum_activity_transport_research.md}
\item \path{work/connected_carrier_coarsening_research.md}
\end{itemize}
\endgroup

\paragraph{Selected audit records.}
\begingroup\small
\begin{itemize}
\item \path{work/coarse_transport_audit.md}
\item \path{work/coarse_stationary_support_provenance_audit.md}
\item \path{work/quantum_activity_transport_theorem_audit.md}
\end{itemize}
\endgroup
Corrections/scope register: C20, C21.\par

\subsection{Notebook 23: gauge covariance of the small-field step}
\label{hnc-23}
\noindent\path{outputs/ym4-gauge-covariance-notebook.tex}\par

\paragraph{Main results.}
\begin{itemize}
\item Exact gauge covariance of the actual small-field chart integral, including stationary centering, independent cutoff, Haar, pivot and Gaussian factors.
\item The complete generated correction obeys the local nonlinear Ward identity; its total linear jet vanishes for a semisimple gauge algebra.
\item An exponential Hessian-row bound and a carefully limited infinite-kernel Fourier consequence are derived.
\end{itemize}

\paragraph{Assumptions and dependencies.}
\begin{itemize}
\item A common local gauge domain and the literal product of invariant color-ball windows; extra indicators must also transform covariantly.
\item Inherited integral identities require convergence; logarithm identities additionally require nonvanishing and compatible analytic branches.
\end{itemize}

Selected upstream notebooks: 19, 20, 21, 22.\par

\paragraph{Limits and later developments.}
\begin{itemize}
\item Finite-torus Ward identities hold at discrete momenta and do not by themselves prove a continuum zero-momentum classification.
\item No separate activity is claimed gauge invariant merely because the total is.
\end{itemize}

Notebook 29 supplies actual spatial symmetries, structural kernel limits and the full low-order nonlinear Ward extraction.

\paragraph{Proof modules and source maps.}
\begingroup\small
\begin{itemize}
\item \path{work/small_field_gauge_covariance_fragment.tex}
\item \path{work/small_field_gauge_covariance_research.md}
\end{itemize}
\endgroup

\paragraph{Selected audit records.}
\begingroup\small
\begin{itemize}
\item \path{work/small_field_gauge_covariance_theorem_audit.md}
\item \path{work/coarse_transport_audit.md}
\end{itemize}
\endgroup
Corrections/scope register: C22.\par

\subsection{Notebook 24: inherited interactions and the next-step norm test}
\label{hnc-24}
\noindent\path{outputs/ym4-inherited-reentry-notebook.tex}\par

\paragraph{Main results.}
\begin{itemize}
\item A directional Cauchy estimate bounds the inherited fluctuation increment by the cutoff displacement divided by the available analytic margin, without a spurious support-size factor.
\item The actual source pivot dependencies enlarge carriers only to specified neighboring blocks.
\item A quantitative logarithm comparison follows under an explicit next-step support-weight inequality.
\end{itemize}

\paragraph{Assumptions and dependencies.}
\begin{itemize}
\item The actual relative fluctuation chart, a strict common input domain and the stated connected input activity norm.
\item The grouped interaction exponent and its multiplier cost are both included; the Gaussian tail threshold is the original linkwise threshold.
\end{itemize}

Selected upstream notebooks: 18, 19, 22, 23.\par

\paragraph{Limits and later developments.}
\begin{itemize}
\item The support-weight inequality is a sufficient condition, not verified merely by declaring the blocking ratio large.
\item The good linkwise block bound does not improve the bad Euclidean tail threshold by the square root of dimension.
\end{itemize}

Notebook 26 verifies a compatible weight scheme through independent auxiliary grouping; Notebook 28 handles a common analytic radius.

\paragraph{Proof modules and source maps.}
\begingroup\small
\begin{itemize}
\item \path{work/inherited_quantum_reentry_fragment.tex}
\item \path{work/inherited_quantum_reentry_research.md}
\end{itemize}
\endgroup

\paragraph{Selected audit records.}
\begingroup\small
\begin{itemize}
\item \path{work/inherited_quantum_reentry_audit.md}
\end{itemize}
\endgroup
Corrections/scope register: C17, C21.\par

\subsection{Notebook 25: a smaller representative of the classical gauge orbit}
\label{hnc-25}
\noindent\path{outputs/ym4-pre-gauge-notebook.tex}\par

\paragraph{Main results.}
\begin{itemize}
\item A pre-root-gauge representative of the classical orbit has a uniformly small link angle proportional to the inverse endpoint ratio.
\item Complete-output-cell Banach Cauchy estimates give a strict supported background family without a false count of every output coordinate.
\item A gauge-invariant total functional analytically extends its original composition germ through this smaller representative.
\end{itemize}

\paragraph{Assumptions and dependencies.}
\begin{itemize}
\item Fixed primitive blocking and actual ordinary prefix averages of the selected classical minimum.
\item The representative is taken before the final root-gauge reset; its gauge relation to the physical branch is explicit.
\end{itemize}

Selected upstream notebooks: 15, 16, 21, 22, 23.\par

\paragraph{Limits and later developments.}
\begin{itemize}
\item The smallness does not hold for the gauge-zero representative on arbitrary pure-gauge data.
\item Analytic continuation of a background total alone does not identify the full inherited fluctuation integral.
\end{itemize}

Notebook 28 constructs the actual joint family needed for the integral. Notebook 29 strengthens the prefix to complete first-gradient Holder support families.

\paragraph{Proof modules and source maps.}
\begingroup\small
\begin{itemize}
\item \path{work/pre_gauge_background_transport_fragment.tex}
\item \path{work/pre_gauge_background_transport_research.md}
\end{itemize}
\endgroup

\paragraph{Selected audit records.}
\begingroup\small
\begin{itemize}
\item \path{work/pre_gauge_background_transport_provenance_audit.md}
\end{itemize}
\endgroup
Corrections/scope register: C21, C23.\par

\subsection{Notebook 26: closing the support-weight compatibility test}
\label{hnc-26}
\noindent\path{outputs/ym4-auxiliary-grouping-notebook.tex}\par

\paragraph{Main results.}
\begin{itemize}
\item Independent auxiliary integration cells separate the physical source blocking ratio from Gaussian support bookkeeping.
\item At fixed physical ratio, sufficiently large auxiliary cells amplify the existing positive decay rate and permit a fixed output carrier exponent.
\item Tunable Gaussian density costs, classical precision reserves and determinant chains close the inherited support-weight compatibility test.
\end{itemize}

\paragraph{Assumptions and dependencies.}
\begin{itemize}
\item The physical ratio is fixed first, auxiliary grouping second, and the common background radius last.
\item The covariance/precision exponents satisfy the strict ordering required by the support algebra; all old and new backgrounds are restricted to their common domain.
\end{itemize}

Selected upstream notebooks: 18, 21, 22, 24, 25.\par

\paragraph{Limits and later developments.}
\begin{itemize}
\item Constants may be very large functions of the chosen geometry.
\item This closes a norm-compatibility test, not repeated-step field-domain or amplitude stability.
\end{itemize}

The composite step and common-domain construction in 27--28 make the resulting integration scheme concrete.

\paragraph{Proof modules and source maps.}
\begingroup\small
\begin{itemize}
\item \path{work/auxiliary_gaussian_grouping_fragment.tex}
\item \path{work/auxiliary_gaussian_grouping_research.md}
\end{itemize}
\endgroup

\paragraph{Selected audit records.}
\begingroup\small
\begin{itemize}
\item \path{work/auxiliary_gaussian_grouping_theorem_audit.md}
\item \path{work/auxiliary_polymer_grouping_audit.md}
\end{itemize}
\endgroup
Corrections/scope register: C24.\par

\subsection{Notebook 27: an exact composite blocking step}
\label{hnc-27}
\noindent\path{outputs/ym4-composite-macrostep-notebook.tex}\par

\paragraph{Main results.}
\begin{itemize}
\item The actual composite ordinary map, rather than a substituted direct-ratio map, has an exact two-block incidence structure and one-consumer pivot.
\item A separate direct-ratio internal contour gauge and a finite-patch nullspace argument give a retained Gaussian uniformly coercive in old depth and volume.
\item An exactly normalized compact gauge insertion and local Haar/coarea formula preserve all constants and the independent cutoff.
\end{itemize}

\paragraph{Assumptions and dependencies.}
\begin{itemize}
\item A fixed primitive ratio and a fixed finite composite endpoint; constants can depend on that endpoint.
\item The selected group mean satisfies the actual source axioms; the source's implicit mean equation is only one admissible realization.
\item A specified measurable equivariant completion is distinguished from its local analytic definition.
\end{itemize}

Selected upstream notebooks: 15, 18, 19, 21, 25, 26.\par

\paragraph{Limits and later developments.}
\begin{itemize}
\item The composite ordinary average is not assumed equal to direct straight averaging away from the flat corrected comparison.
\item Local small-field integral identities do not yet bound the rejected configurations.
\end{itemize}

Notebook 28 integrates this exact composite family on one common analytic ball.

\paragraph{Proof modules and source maps.}
\begingroup\small
\begin{itemize}
\item \path{work/composite_macrostep_gaussian_fragment.tex}
\item \path{work/composite_macrostep_geometry_fragment.tex}
\item \path{work/composite_macrostep_integral_fragment.tex}
\item \path{work/composite_macrostep_gaussian_research.md}
\item \path{work/composite_macrostep_geometry_research.md}
\item \path{work/composite_macrostep_integral_research.md}
\end{itemize}
\endgroup

\paragraph{Selected audit records.}
\begingroup\small
\begin{itemize}
\item \path{work/composite_macrostep_gaussian_audit.md}
\item \path{work/composite_macrostep_integral_rotation_audit.md}
\item \path{work/composite_macrostep_integral_theorem_audit.md}
\end{itemize}
\endgroup
Corrections/scope register: C25.\par

\subsection{Notebook 28: a fixed analytic ball and high-degree stability}
\label{hnc-28}
\noindent\path{outputs/ym4-common-domain-notebook.tex}\par

\paragraph{Main results.}
\begin{itemize}
\item A local-core/small-tail expansion composes the actual pre-gauge joint background/fluctuation family on one fixed analytic ball.
\item Classical precision and measure factors are substituted before Gaussian integration, so the output is already in the required coarse variable.
\item Coefficientwise degree-five-and-higher background terms contract by an inverse endpoint ratio; the actual quantum response to differences with matching low jets adds only a small inherited fluctuation error.
\end{itemize}

\paragraph{Assumptions and dependencies.}
\begin{itemize}
\item The ordered choices of endpoint, auxiliary grouping and final radius; fixed strict tail and precision budgets.
\item A shared canonical core/tail grouping identifies the background part of the joint integral coefficientwise.
\item The high-degree result assumes each input coefficient has zero Taylor jets through degree four.
\end{itemize}

Selected upstream notebooks: 18, 21, 25, 26, 27.\par

\paragraph{Limits and later developments.}
\begin{itemize}
\item Coefficientwise Taylor projection is not gauge invariant and its high-degree subspace is not asserted invariant under the quantum map.
\item A common domain and a difference contraction do not alone control all regenerated low jets or amplitudes.
\end{itemize}

Notebook 29 identifies the total low-order gauge structure; Notebooks 30--31 supply quotient extraction and finite amplitude induction.

\paragraph{Proof modules and source maps.}
\begingroup\small
\begin{itemize}
\item \path{work/common_domain_classical_support_fragment.tex}
\item \path{work/common_domain_joint_transport_fragment.tex}
\item \path{work/common_domain_macrostep_fragment.tex}
\item \path{work/high_jet_background_transport_fragment.tex}
\item \path{work/high_jet_quantum_difference_fragment.tex}
\item \path{work/stationary_before_gaussian_research.md}
\item \path{work/composite_macrostep_joint_family_research.md}
\item \path{work/high_jet_background_transport_research.md}
\item \path{work/high_jet_quantum_difference_research.md}
\end{itemize}
\endgroup

\paragraph{Selected audit records.}
\begingroup\small
\begin{itemize}
\item \path{work/common_domain_classical_support_theorem_audit.md}
\item \path{work/composite_common_domain_fragment_provenance_audit.md}
\item \path{work/high_jet_background_transport_theorem_audit.md}
\item \path{work/high_jet_quantum_difference_rotation_audit.md}
\end{itemize}
\endgroup

Checkpoint: \path{work/after_notebook28_state.txt}.\par
Corrections/scope register: C23, C26.\par

\subsection{Notebook 29: total Ward extraction and classical spatial gain}
\label{hnc-29}
\noindent\path{outputs/ym4-gauge-extraction-notebook.tex}\par

\paragraph{Main results.}
\begin{itemize}
\item Actual spatial and gauge symmetries, exponential kernels and joint structural limits identify the normalized Maxwell quadratic coefficient.
\item Nonlinear Ward identities determine the compatible cubic and quartic low-momentum subtraction by a full gauge-invariant classical reference.
\item Complete Holder prefix families and an exact two-difference Laurent factorization yield a connected-norm spatial gain; the classical semigroup permits the final endpoint to be chosen after auxiliary grouping.
\item Analytic interaction compactness and a Banach quotient by equality of normalized totals make representation changes and bounded reference extraction precise.
\end{itemize}

\paragraph{Assumptions and dependencies.}
\begin{itemize}
\item Coherent infinite structural systems or exact periodizations of their kernels; full hypercubic symmetry uses cofinal cubic tori.
\item First-gradient Holder estimates are for complete fibers with consistent support assignments.
\item The endpoint used for contraction is chosen after the fixed grouping constants.
\end{itemize}

Selected upstream notebooks: 07, 15, 16, 21, 23, 25, 28.\par

\paragraph{Limits and later developments.}
\begin{itemize}
\item Arbitrary finite winding kernels are not identified with infinite periodizations.
\item The spatial constant can depend on auxiliary grouping; a nominal inverse-ratio power is not a contraction if parameters are chosen in the wrong order.
\item The naive symmetric quadratic difference identity is false for odd-skew matrix kernels and is not used.
\end{itemize}

Notebook 30 closes the universal gauge/quantum quotient interfaces; Notebook 31 gives the finite-regulator repair.

\paragraph{Proof modules and source maps.}
\begingroup\small
\begin{itemize}
\item \path{work/activity_quotient_fragment.tex}
\item \path{work/analytic_interaction_compactness_fragment.tex}
\item \path{work/classical_background_semigroup_fragment.tex}
\item \path{work/low_kernel_spatial_transport_fragment.tex}
\item \path{work/macrostep_symmetry_fragment.tex}
\item \path{work/pregauge_holder_prefix_fragment.tex}
\item \path{work/total_kernel_ward_extraction_fragment.tex}
\item \path{work/composite_macrostep_symmetry_research.md}
\item \path{work/low_kernel_spatial_transport_research.md}
\item \path{work/classical_background_semigroup_research.md}
\item \path{work/analytic_interaction_compactness_research.md}
\item \path{work/activity_quotient_research.md}
\end{itemize}
\endgroup

\paragraph{Selected audit records.}
\begingroup\small
\begin{itemize}
\item \path{work/composite_macrostep_symmetry_theorem_audit.md}
\item \path{work/total_kernel_ward_extraction_rotation_audit.md}
\item \path{work/low_kernel_spatial_transport_rotation_audit.md}
\item \path{work/classical_background_semigroup_rotation_audit.md}
\item \path{work/activity_quotient_theorem_audit.md}
\end{itemize}
\endgroup

Checkpoint: \path{work/after_notebook29_state.txt}.\par
Corrections/scope register: C22, C27, C28.\par

\subsection{Notebook 30: quantum descent and renormalized block stability}
\label{hnc-30}
\noindent\path{outputs/ym4-quantum-descent-notebook.tex}\par

\paragraph{Main results.}
\begin{itemize}
\item A summable connecting gauge, compact approximation and domain margin give a universal classical action on invariant quotient classes.
\item Exact finite periodization and strict Gaussian local continuity identify the actual quantum map and compatible finite histories on that quotient.
\item Transport of the actual reference makes repeated extraction telescope. A finite block of quantum steps has a stable O(b\textasciicircum{}2) amplitude envelope while couplings stay weak.
\item A fixed outer analytic radius provides the matching quadratic zero-input force.
\end{itemize}

\paragraph{Assumptions and dependencies.}
\begin{itemize}
\item A jointly compatible homogeneous structural system, the fixed local quantum budgets and allowed compact gauge paths.
\item Finite-history identification requires uniformly bounded actual intermediate representatives; continuity uses carrier-size, order and geometric tails.
\item The weak coupling sequence is prescribed or known to remain inside the admissible range.
\end{itemize}

Selected upstream notebooks: 28, 29.\par

\paragraph{Limits and later developments.}
\begin{itemize}
\item This does not itself establish actual running couplings or regulator-uniform representatives at arbitrarily many finite steps.
\item The zero-input force is not O(g\textasciicircum{}2) as a whole: pivot and Gaussian determinants persist at zero coupling.
\end{itemize}

Notebook 31 closes the finite representative-budget interface; Notebook 37 supplies the positive drift needed for tuning. The generic adapted marginal coordinate remains a fallback rather than the actual final formulation.

\paragraph{Proof modules and source maps.}
\begingroup\small
\begin{itemize}
\item \path{work/classical_summable_gauge_fragment.tex}
\item \path{work/extracted_macrostep_recurrence_fragment.tex}
\item \path{work/periodic_interaction_descent_fragment.tex}
\item \path{work/quantum_local_limit_fragment.tex}
\item \path{work/zero_input_quadratic_force_fragment.tex}
\item \path{work/classical_summable_gauge_research.md}
\item \path{work/quantum_local_limit_research.md}
\item \path{work/extracted_macrostep_recurrence_rotation_research.md}
\item \path{work/quantum_zero_input_force_research.md}
\end{itemize}
\endgroup

\paragraph{Selected audit records.}
\begingroup\small
\begin{itemize}
\item \path{work/classical_summable_gauge_theorem_audit.md}
\item \path{work/periodic_interaction_descent_theorem_audit.md}
\item \path{work/quantum_local_limit_provenance_audit.md}
\item \path{work/extracted_macrostep_recurrence_provenance_audit.md}
\item \path{work/quantum_zero_input_force_rotation_audit.md}
\end{itemize}
\endgroup

Checkpoint: \path{work/after_notebook30_state.txt}.\par
Corrections/scope register: C29, C30.\par

\subsection{Notebook 31: finite local induction and coupling extraction}
\label{hnc-31}
\noindent\path{outputs/ym4-finite-induction-notebook.tex}\par

\paragraph{Main results.}
\begin{itemize}
\item Finite Ward repair and direct periodic localization retain wrapping errors at the same field radius and carrier exponent.
\item Exact finite quotient descent and equivalent small representatives give local induction through every permitted finite history; vacuum compensation retains the original raw density.
\item Fixed-threshold parity isolates the leading zero-input pivot/Gaussian determinant.
\item A quadratic-only fixed-radius semigroup argument proves exact classical Maxwell preservation in the compatible infinite system.
\end{itemize}

\paragraph{Assumptions and dependencies.}
\begin{itemize}
\item Real extracted local trajectories whose couplings remain uniformly weak and whose periods remain above the stated threshold.
\item Exact representative replacement is equality of normalized totals, not accumulation of approximation errors.
\item The infinite Maxwell-preservation proof does not use a shrinking high-jet domain.
\end{itemize}

Selected upstream notebooks: 28, 29, 30.\par

\paragraph{Limits and later developments.}
\begin{itemize}
\item Finite cutoff extraction is not exactly preserved; its leakage must still be controlled.
\item The theorem concerns the homogeneous admitted local sector, not the full Wilson history or arbitrary effective large-field measures.
\end{itemize}

Notebook 32 controls finite leakage qualitatively and formulates shooting; Notebook 37 verifies its positive leading margin.

\paragraph{Proof modules and source maps.}
\begingroup\small
\begin{itemize}
\item \path{work/classical_marginal_transport_fragment.tex}
\item \path{work/finite_classical_contraction_fragment.tex}
\item \path{work/finite_quantum_induction_fragment.tex}
\item \path{work/quantum_zero_input_parity_fragment.tex}
\item \path{work/finite_classical_contraction_research.md}
\item \path{work/finite_quantum_induction_research.md}
\item \path{work/classical_marginal_transport_research.md}
\end{itemize}
\endgroup

\paragraph{Selected audit records.}
\begingroup\small
\begin{itemize}
\item \path{work/finite_classical_contraction_rotation_audit.md}
\item \path{work/finite_quantum_induction_rotation_audit.md}
\item \path{work/quantum_zero_input_parity_rotation_audit.md}
\item \path{work/classical_marginal_transport_rotation_audit.md}
\end{itemize}
\endgroup

Checkpoint: \path{work/after_notebook31_state.txt}.\par
Corrections/scope register: C30, C31.\par

\subsection{Notebook 32: determinant composition and finite coupling tuning}
\label{hnc-32}
\noindent\path{outputs/ym4-determinant-flow-notebook.tex}\par

\paragraph{Main results.}
\begin{itemize}
\item Finite classical marginal leakage vanishes uniformly over varying old depths along cofinal nested periods; the finite quantum map is continuous in its class and positive coupling.
\item The actual leading determinant has a convergent Hessian trace with both chart-curvature contributions and a finite Haar/Laplace composition cocycle.
\item An exact gauge/coarea bridge retains vector, based ghost and coarse constraint factors; the shifted coarse effective determinant has a uniform Maxwell bound.
\item Uniform positivity over all structural limits transfers qualitatively to actual finite blocks, and a separate shooting theorem tunes long finite histories under that sign hypothesis.
\end{itemize}

\paragraph{Assumptions and dependencies.}
\begin{itemize}
\item Fixed finite block length, cofinal compatible periods, actual reference normalization and local minima on a common real neighborhood.
\item Joint extraction includes every shifted old depth; the positive margin must cover all obtainable structural limits, not one favorable subsequence.
\item Stationary phase is finite dimensional on the correctly normalized orbit/quotient tube.
\end{itemize}

Selected upstream notebooks: 15, 16, 21, 29, 30, 31.\par

\paragraph{Limits and later developments.}
\begin{itemize}
\item There is no quantitative finite-volume convergence rate or summable-rate claim.
\item The leading cocycle does not assert composition of complete positive-coupling admitted sectors.
\item The based ghost has no uniform block-scale mass gap.
\end{itemize}

Notebooks 33--37 evaluate the actual leading sign. The canonical coarse-completion simplification is separate from the fixed positive shift used for uniform bounds.

\paragraph{Proof modules and source maps.}
\begingroup\small
\begin{itemize}
\item \path{work/classical_one_loop_cocycle_fragment.tex}
\item \path{work/determinant_gauge_bridge_fragment.tex}
\item \path{work/determinant_maxwell_trace_fragment.tex}
\item \path{work/finite_block_drift_transfer_fragment.tex}
\item \path{work/finite_classical_marginal_leakage_fragment.tex}
\item \path{work/finite_quantum_coupling_continuity_fragment.tex}
\item \path{work/finite_trajectory_shooting_fragment.tex}
\item \path{work/determinant_maxwell_trace_research.md}
\item \path{work/classical_one_loop_cocycle_research.md}
\item \path{work/determinant_gauge_bridge_research.md}
\item \path{work/finite_leading_drift_transfer_research.md}
\end{itemize}
\endgroup

\paragraph{Selected audit records.}
\begingroup\small
\begin{itemize}
\item \path{work/finite_classical_marginal_leakage_rotation_audit.md}
\item \path{work/finite_quantum_coupling_continuity_rotation_audit.md}
\item \path{work/determinant_maxwell_trace_theorem_audit.md}
\item \path{work/classical_one_loop_cocycle_rotation_audit.md}
\item \path{work/determinant_gauge_bridge_rotation_audit.md}
\item \path{work/finite_leading_drift_transfer_provenance_audit.md}
\item \path{work/finite_trajectory_shooting_rotation_audit.md}
\end{itemize}
\endgroup

Checkpoint: \path{work/after_notebook32_state.txt}.\par
Corrections/scope register: C31, C32, C33.\par

\subsection{Notebook 33: gauge-invariant determinant densities and Cartan probes}
\label{hnc-33}
\noindent\path{outputs/ym4-cartan-trace-notebook.tex}\par

\paragraph{Main results.}
\begin{itemize}
\item Individually gauge-invariant diagonal determinant densities are related to the total Maxwell moment.
\item Actual stationary affine Cartan jets remove the required multiplier-curvature terms, including for effective classical actions.
\item An exactly normalized massive continuum reference and an actual heat identity isolate the comparison problem; the based-ghost ultraviolet estimate is established.
\end{itemize}

\paragraph{Assumptions and dependencies.}
\begin{itemize}
\item The chosen joint infinite structural system, actual stationary jet identities and common determinant conventions.
\item The reference computation and actual lattice comparison are kept distinct, including signs and real color multiplicities.
\end{itemize}

Selected upstream notebooks: 07, 16, 21, 29, 32.\par

\paragraph{Limits and later developments.}
\begin{itemize}
\item The massive continuum logarithm alone did not establish the actual lattice coefficient.
\item Polynomial affine probes are source jets in a local structural system, not an assumed finite-torus globally periodic flux family.
\end{itemize}

Notebook 34 computes the whole based ghost; 35--37 supply the joint infrared and actual vector comparison.

\paragraph{Proof modules and source maps.}
\begingroup\small
\begin{itemize}
\item \path{work/actual_heat_comparison_fragment.tex}
\item \path{work/cartan_determinant_moment_fragment.tex}
\item \path{work/gauge_invariant_determinant_density_fragment.tex}
\item \path{work/massive_cartan_reference_fragment.tex}
\item \path{work/mean_zero_one_loop_source_interface.txt}
\end{itemize}
\endgroup

\paragraph{Selected audit records.}
\begingroup\small
\begin{itemize}
\item \path{work/cartan_trace_root_audit.txt}
\item \path{work/cartan_trace_provenance_audit.txt}
\item \path{work/cartan_heat_rotation_audit.txt}
\end{itemize}
\endgroup

Checkpoint: \path{work/after_notebook33_state.txt}.\par
Corrections/scope register: C33, C34.\par

\subsection{Notebook 34: ghost asymptotics and the physical determinant}
\label{hnc-34}
\noindent\path{outputs/ym4-pinned-ghost-notebook.tex}\par

\paragraph{Main results.}
\begin{itemize}
\item A returning-walk area formula and Bessel identities give the exact unpinned ghost heat coefficient and controlled pinned prefix.
\item The point-pinned gap and fourth-area-moment estimate show that the complete based-ghost coefficient has an extra logarithm from times K\textasciicircum{}2 to K\textasciicircum{}4.
\item The actual full composite vector ultraviolet term is bounded without assuming the false K\textasciicircum{}-1 norm for the ordinary constraint.
\item An exact physical-quotient identity, projected constraint norm and flat physical gap separate genuine cancellation from contour-gauge artifacts.
\end{itemize}

\paragraph{Assumptions and dependencies.}
\begin{itemize}
\item The actual based scalar, ordinary composite constraint and correctly normalized finite or structural-system setting of each theorem.
\item The projected estimate applies to the transverse coarse projection, not the full ordinary constraint.
\end{itemize}

Selected upstream notebooks: 07, 08, 27, 32, 33.\par

\paragraph{Limits and later developments.}
\begin{itemize}
\item A scalar ghost tail bounded after K\textasciicircum{}2 is false; the full vector must cancel its extra logarithm.
\item The unprojected contour-gauge matrix has a counterexample to the desired K\textasciicircum{}-2 gap.
\item Separately regularized physical/scalar factors are not yet the actual determinant.
\end{itemize}

Notebook 35 replaces point pinning by covariant zero block means and treats the joint infrared continuation.

\paragraph{Proof modules and source maps.}
\begingroup\small
\begin{itemize}
\item \path{work/physical_quotient_determinant_fragment.tex}
\item \path{work/pinned_ghost_long_time_fragment.tex}
\item \path{work/pinned_ghost_prefix_fragment.tex}
\item \path{work/vector_short_time_fragment.tex}
\item \path{work/contour_heat_comparison_research.txt}
\item \path{work/vector_short_time_research.txt}
\end{itemize}
\endgroup

\paragraph{Selected audit records.}
\begingroup\small
\begin{itemize}
\item \path{work/pinned_ghost_prefix_provenance_audit.txt}
\item \path{work/pinned_ghost_long_time_provenance_audit.txt}
\item \path{work/vector_short_time_provenance_audit.txt}
\item \path{work/physical_quotient_gap_rotation_audit.txt}
\item \path{work/contour_heat_comparison_provenance_audit.txt}
\end{itemize}
\endgroup

Checkpoint: \path{work/after_notebook34_state.txt}.\par
Corrections/scope register: C33, C35.\par

\subsection{Notebook 35: infrared cancellation and physical heat bounds}
\label{hnc-35}
\noindent\path{outputs/ym4-infrared-cancellation-notebook.tex}\par

\paragraph{Main results.}
\begin{itemize}
\item A constant-holonomy example demonstrates noncommuting massless and thermodynamic scalar limits and the required physical compensation.
\item Moving-reference covariant scalar averages have an exact constant row Gram matrix; a joint removable determinant identity replaces pinning by zero covariant block means.
\item Scalar and vector factors have flat K\textasciicircum{}-2 gaps. Block-mean and tent lifting give Laplacian domination, heat and inverse-power bounds.
\item Local diagonal densities identify the actual fixed-endpoint Maxwell coefficient without singular factorwise logarithms at flatness.
\end{itemize}

\paragraph{Assumptions and dependencies.}
\begin{itemize}
\item The actual endpoint tangent identities and joint continuation from the irreducible local germ.
\item Counting versus physical normalization, harmonic modes and the scalar complement are kept explicit.
\end{itemize}

Selected upstream notebooks: 07, 08, 27, 32, 33, 34.\par

\paragraph{Limits and later developments.}
\begin{itemize}
\item A flat gap or Schatten bound does not imply uniformly small raw curvature vertices.
\item The hard projector admits a pure-gauge relative-vertex counterexample; independent scalar/vector regularizations remain invalid.
\end{itemize}

Notebook 36 uses cancellation-preserving soft factors and relative forms to obtain the required large-endpoint comparisons.

\paragraph{Proof modules and source maps.}
\begingroup\small
\begin{itemize}
\item \path{work/covariant_corrected_scalar_fragment.tex}
\item \path{work/covariant_mean_zero_density_fragment.tex}
\item \path{work/joint_scalar_infrared_fragment.tex}
\item \path{work/mean_zero_gauge_spectral_fragment.tex}
\item \path{work/physical_quotient_localization_fragment.tex}
\item \path{work/physical_vector_spectral_fragment.tex}
\item \path{work/physical_quotient_localization_research.txt}
\item \path{work/covariant_corrected_scalar_research.txt}
\item \path{work/mean_zero_gauge_spectral_research.txt}
\item \path{work/physical_vector_curvature_research.txt}
\end{itemize}
\endgroup

\paragraph{Selected audit records.}
\begingroup\small
\begin{itemize}
\item \path{work/joint_scalar_infrared_rotation_audit.txt}
\item \path{work/mean_zero_gauge_bridge_rotation_audit.txt}
\item \path{work/covariant_mean_zero_provenance_audit.txt}
\item \path{work/physical_vector_spectral_provenance_audit.txt}
\item \path{work/mean_zero_gauge_vertex_provenance_audit.txt}
\end{itemize}
\endgroup

Checkpoint: \path{work/after_notebook35_state.txt}.\par
Corrections/scope register: C35, C36.\par

\subsection{Notebook 36: soft determinants and bounded coarse corrections}
\label{hnc-36}
\noindent\path{outputs/ym4-soft-determinant-notebook.tex}\par

\paragraph{Main results.}
\begin{itemize}
\item The actual determinant is factored into a soft fine scalar/vector pair and two coarse determinant corrections with uniformly bounded Maxwell contributions.
\item A corrected ordered-Laurent and local gauge argument controls relative vertices of the actual effective Hessian, with rooted source bounds.
\item The current ultraviolet interval and the physical infrared tail of the soft pair are uniformly bounded; the latter retains the local block projector in a cubic-resolvent Schatten estimate.
\end{itemize}

\paragraph{Assumptions and dependencies.}
\begin{itemize}
\item Actual infinite structural kernels with the proved continuous moments, or their explicit periodizations; the on-shell affine matching is used only where established.
\item Source localization carries a full tree through input/output and source roots, with physical fine-site Lipschitz weights.
\end{itemize}

Selected upstream notebooks: 21, 29, 33, 34, 35.\par

\paragraph{Limits and later developments.}
\begin{itemize}
\item The componentwise-even difference-square shortcut is false and has been replaced.
\item These results reduce the logarithm to the intermediate band; spectral bounds alone do not determine its coefficient.
\item Arbitrary finite winding kernels are not silently treated as exact infinite periodizations.
\end{itemize}

Notebook 37 closes the intermediate physical band and compares the actual base to the Wilson and continuum references.

\paragraph{Proof modules and source maps.}
\begingroup\small
\begin{itemize}
\item \path{work/covariant_hessian_relative_vertices_fragment.tex}
\item \path{work/soft_pair_infrared_tail_fragment.tex}
\item \path{work/soft_pair_short_time_fragment.tex}
\item \path{work/soft_scalar_coarse_reduction_fragment.tex}
\item \path{work/soft_vector_coarse_reduction_fragment.tex}
\item \path{work/covariant_hessian_relative_vertices_research.txt}
\item \path{work/soft_vector_coarse_reduction_research.txt}
\item \path{work/soft_pair_infrared_tail_research.txt}
\end{itemize}
\endgroup

\paragraph{Selected audit records.}
\begingroup\small
\begin{itemize}
\item \path{work/soft_determinant_root_audit.txt}
\item \path{work/covariant_hessian_relative_vertices_rotation_audit.txt}
\item \path{work/soft_vector_coarse_reduction_rotation_audit.txt}
\item \path{work/soft_pair_infrared_tail_rotation_audit.txt}
\item \path{work/soft_pair_short_time_rotation_audit.txt}
\end{itemize}
\endgroup

Checkpoint: \path{work/after_notebook36_state.txt}.\par
Corrections/scope register: C36, C37.\par

\subsection{Notebook 37: the actual logarithm and finite local ultraviolet tuning}
\label{hnc-37}
\noindent\path{outputs/ym4-positive-drift/main.tex}\par

\paragraph{Main results.}
\begin{itemize}
\item Source-localized sectorial resolvents remove the soft penalties on the entire remaining physical band with an integrable O(u) heat error.
\item The nonlinear Ward jet orders (4,3,2), including left-relative chart-current terms, compare the actual base to the Wilson base with O(s\textasciicircum{}-1) error.
\item A direct Wilson-to-continuum rescaling and transverse diagram calculation give O(s\textasciicircum{}-1/2) error and the exact normalization beta(D\_M,K)=22 alpha\_g log K+O(1), uniformly in old depth and every applicable structural system.
\item The resulting positive finite block margin permits arbitrarily long tuned weak local histories, with bare coupling of inverse-square-root order in length.
\end{itemize}

\paragraph{Assumptions and dependencies.}
\begin{itemize}
\item All explicit classical structural, local quantum, determinant and compactness inputs of the preceding series.
\item Uniformity is over every required joint structural limit, not a chosen favorable one; the fixed local step parameters are chosen in the stated order.
\end{itemize}

Selected upstream notebooks: 29, 31, 32, 33, 34, 35, 36.\par

\paragraph{Limits and later developments.}
\begin{itemize}
\item This closes the local leading-drift/tuning problem, not the full Wilson large-field or mixed-history problem.
\item Internal proof audits and compilation are not external mathematical certification.
\item Continuum contact terms are handled by an explicit diagram calculation, not an unjustified absolute spatial-moment assertion.
\end{itemize}

Notebook 38 begins the genuinely normalized full-measure interface and records the remaining global obligations.

\paragraph{Proof modules and source maps.}
\begingroup\small
\begin{itemize}
\item \path{work/actual_base_ward_heat_comparison_fragment.tex}
\item \path{work/actual_positive_beta_assembly_fragment.tex}
\item \path{work/local_positive_drift_fragment.tex}
\item \path{work/soft_penalty_physical_band_fragment.tex}
\item \path{work/soft_scalar_complete_coefficient_fragment.tex}
\item \path{work/wilson_continuum_heat_comparison_fragment.tex}
\item \path{work/actual_base_ward_heat_comparison_research.txt}
\item \path{work/wilson_continuum_heat_comparison_research.txt}
\item \path{work/soft_penalty_physical_band_complete_research.txt}
\end{itemize}
\endgroup

\paragraph{Selected audit records.}
\begingroup\small
\begin{itemize}
\item \path{work/actual_base_ward_heat_comparison_root_audit.txt}
\item \path{work/actual_base_ward_heat_comparison_theorem_audit.txt}
\item \path{work/wilson_continuum_heat_comparison_theorem_audit.txt}
\item \path{work/soft_penalty_physical_band_root_audit.txt}
\item \path{work/local_positive_drift_theorem_audit.txt}
\end{itemize}
\endgroup

Checkpoint: \path{work/after_notebook37_state.txt}.\par
Corrections/scope register: C01, C34, C38.\par

\subsection{Notebook 38: positive full-measure embeddings and local defect bounds}
\label{hnc-38}
\noindent\path{outputs/ym4-full-measure-bridge/main.tex}\par

\paragraph{Main results.}
\begin{itemize}
\item A specified global Borel completion embeds the admitted local sector as a positive part of the actual finite Wilson pushforward, with a normalized output-orbit acceptance and exact positive remainder histories.
\item A full-normalization reflection-positive Peierls estimate controls prescribed bad plaquettes/cells and cluster counts.
\item An actual fixed-box conditional Wilson Laplace theorem and a boundary-only gauge construction give unconditional separated-box defect estimates.
\item A curvature-oscillation counterexample and primary-source comparison identify the limitations of replacing a complete large-field analysis by plaquette smallness.
\end{itemize}

\paragraph{Assumptions and dependencies.}
\begin{itemize}
\item Positive bare Wilson measure, normalized Haar and compatible reflection tilings; fixed embedded boxes with disjoint vertex sets for the box bridge.
\item The literal sector uses invariant color balls and actual gauge-covariant charts. Its local-germ identity needs a common finite-volume orbit domain and containment of every retained configuration.
\item The chosen global completion is part of the construction, not an identification with an unspecified original-source completion.
\end{itemize}

Selected upstream notebooks: 19, 27, 28, 31, 37.\par

\paragraph{Limits and later developments.}
\begin{itemize}
\item The positive remainder has no proved complete mixed-history polymer bound.
\item The fixed-box minimum is an unconstrained conditional bare minimum, not the nonlinear coarse-fiber RG minimum.
\item Small plaquettes do not yield uniform physical C1,alpha control; arbitrary growing components, holonomy sectors and later effective measures require additional work.
\end{itemize}

Notebook 39 adds the final normalized fixed-box replacement theorem. It preserves the exterior and partition mass and gives local comparison bounds, but does not remove the RG constrained-window interface.

\paragraph{Proof modules and source maps.}
\begingroup\small
\begin{itemize}
\item \path{work/fixed_box_wilson_tail_fragment.tex}
\item \path{work/full_measure_box_bridge_fragment.tex}
\item \path{work/full_measure_geometric_interface_fragment.tex}
\item \path{work/full_wilson_peierls_fragment.tex}
\item \path{work/full_wilson_positive_split_fragment.tex}
\item \path{work/positive_finite_measure_embedding_research.txt}
\item \path{work/full_wilson_large_field_peierls_research.txt}
\item \path{work/full_measure_large_field_interface_research.txt}
\end{itemize}
\endgroup

\paragraph{Selected audit records.}
\begingroup\small
\begin{itemize}
\item \path{work/full_measure_bridge_root_audit.txt}
\item \path{work/positive_finite_measure_embedding_rotation_audit.txt}
\item \path{work/full_wilson_peierls_theorem_audit.txt}
\item \path{work/fixed_box_wilson_tail_provenance_audit.txt}
\item \path{work/full_measure_box_bridge_provenance_audit.txt}
\item \path{work/full_measure_geometric_interface_theorem_audit.txt}
\end{itemize}
\endgroup

Checkpoint: \path{work/after_notebook38_state.txt}.\par
Corrections/scope register: C19, C39, C40.\par

\subsection{Notebook 39: normalized local replacement in the bare Wilson measure}
\label{hnc-39}
\noindent\path{outputs/ym4-normalized-box-replacement/main.tex}\par

\paragraph{Main results.}
\begin{itemize}
\item An exactly normalized local replacement on arbitrary finite collections of vertex-disjoint fixed bare-Wilson boxes preserves the entire exterior marginal, gauge invariance and raw partition mass.
\item Good-boundary conditional interiors are replaced by their actual conditional window laws; bad-boundary interiors are unchanged. A positive failure expansion and explicit coupling give iid domination of discrepancy marks.
\item Local total variation, entropy, bounded-source ratios, connected cumulants and cluster bounds are uniform in the number of boxes.
\end{itemize}

\paragraph{Assumptions and dependencies.}
\begin{itemize}
\item The exact conditional factorization and actual normalized fixed-box tail/Peierls inputs of Notebook 38.
\item Fixed embedded box geometry, compatible periodic surroundings, T0<=T<=r/g and a=C exp(-c T\textasciicircum{}2)<=1/2.
\item Windows are lifted from the fixed-box forest coordinates with invariant edge/color norms; all denominators are exact conditional masses.
\end{itemize}

Selected upstream notebooks: 38.\par

\paragraph{Limits and later developments.}
\begin{itemize}
\item The independent domination applies to discrepancy marks, not to correlated bad-boundary events.
\item The TV equality includes the entire exterior; a marginal forgetting it has only an upper bound. Global TV may approach one.
\item This is not the RG coarse-fiber window, a growing-component replacement, a later effective-action theorem or a renormalized source hierarchy.
\end{itemize}

This is the final stopping theorem of the handoff. The full mixed-history, continuum reconstruction and Yang--Mills identification tasks remain separate.

\paragraph{Proof modules and source maps.}
\begingroup\small
\begin{itemize}
\item \path{work/normalized_box_replacement_fragment.tex}
\item \path{work/normalized_box_replacement_research.txt}
\end{itemize}
\endgroup

\paragraph{Selected audit records.}
\begingroup\small
\begin{itemize}
\item \path{work/normalized_box_replacement_rotation_audit.txt}
\end{itemize}
\endgroup
Corrections/scope register: C40, C41.\par

\subsection{How to preserve and reuse the archive}
\label{hnc-archive-policy}
Preserve the authoritative outputs, their PDFs, all listed fragments,
research notes, audits, checkpoints, builders and source provenance.
Keep unlisted working drafts as historical evidence, with a clear
non-authoritative label; deleting a failed argument would lose useful
information about why the current proof has its present hypotheses.
Do not treat a chronological filename or an old audit pass as permission
to ignore a later correction.

Record the source checkout's pinned revision and retrieval instructions
separately from the new work. Include primary URLs and extraction notes, and explicitly
flag missing or zero-byte local PDF downloads. Compilation and image
review logs support reproducibility and readability, not the truth of
the mathematical claims. Exclude recursive copies of the handoff
directory itself. The source paths in this catalogue can be mapped
directly into the portable bundle by its packager.

